\documentclass[11pt]{article}

\usepackage[margin=1in]{geometry}
\usepackage{amsmath,amssymb,booktabs,tabularx}

\title{Terminal cumulant adapter}
\author{}
\date{}

\begin{document}
\maketitle

\noindent\textbf{Status.} Draft\\
\textbf{Plan item.} Terminal cumulant adapter
\((v,\rho_\Omega\mapsto u_\pm^B,u_\pm^E,z_\pm,r_\Omega,\rho_\Omega,q)\)\\
\textbf{Depends on.} \texttt{03\_conditional\_cf\_mgf\_trading\_block.tex},
\texttt{04\_affine\_transform\_ansatz.tex},
\texttt{05\_riccati\_ode\_system\_for\_c\_d\_v.tex},
\texttt{06\_affine\_domain\_constraints.tex}\\
\textbf{Feeds.} Pricing density normalization, martingale tilt, risk-neutral characteristic function

\section{Purpose}

Derive the adapter that turns a terminal log-moment argument into the affine
schedule query parameters.
This is the bridge between the terminal cumulant \(K_{t,T}\) used by pricing and
the affine oracle's activity, event-count, market-pressure anchor, and
news-return loadings.

\section{Setup}

Let \(t\) be valuation time and \(T\) the terminal option expiry.
The observed log return is decomposed into a deterministic terminal component,
market-pressure anchor movement, and mean-reversion-weighted Janiform core
shocks:
\[
    R_{t,T}
    =
    d_{t,T}
    +
    A_{t,T}^{\Omega}
    +
    \sum_{h=t}^{T-1} w_h X_h^{\mathrm{core}},
\]
where
\[
    A_{t,T}^{\Omega}
    =
    \gamma_\Omega B_{t,T}^{\Omega},
    \qquad
    B_{t,T}^{\Omega}
    =
    \int_t^T b_{\Omega_s}\,dE_s^\Omega.
\]
The weight \(w_h\ge0\) is the terminal mean-reversion weight applied to interval
\(h\)'s core shock.
In the no-observed-state-mean-reversion case, \(w_h=1\).
The market-pressure anchor-pressure movement \(A_{t,T}^{\Omega}\) is
not multiplied by \(w_h\).

The terminal stochastic cumulant is
\[
    K_{t,T}(v,\gamma_+,\gamma_-,\rho_\Omega)
    =
    \log\mathbb E^P\!\left[
        \exp\!\left(
            vA_{t,T}^{\Omega}
            +v\sum_h w_h X_h^{\mathrm{core}}
            +\gamma_+N_{t,T}^+
            +\gamma_-N_{t,T}^-
            -\rho_\Omega B_{t,T}^{\Omega}
        \right)
        \mid\mathcal F_t
    \right].
\]
The deterministic component \(d_{t,T}\), including the known contribution from
the initial residual component, is kept outside \(K_{t,T}\).

\section{Notation}

\begin{tabularx}{\textwidth}{lXX}
\toprule
Symbol & Meaning & Units or domain\\
\midrule
\(v\) & Terminal log-moment argument & Complex transform domain\\
\(\tilde v_h\) & Weighted interval argument, \(\tilde v_h=w_hv\) & Complex transform domain\\
\(w_h\) & Terminal return weight for interval \(h\) & Non-negative finite\\
\(A_{t,T}^{\Omega}\) & Market-pressure anchor-pressure movement & Stochastic finite log return\\
\(B_{t,T}^{\Omega}\) & Signed-balance exposure integral & Real finite\\
\(r_{\Omega,h}\) & Terminal return loading passed to pressure transform & Complex finite\\
\(\gamma_\Omega\) & Anchor drift per signed-balance exposure & Real finite\\
\(\rho_\Omega\) & Signed-imbalance exposure risk price & Real finite\\
\(e_h^\Omega\) & Interval signed-balance exposure scale & Non-negative finite\\
\(\gamma_\pm\) & Log activity premia, \(\gamma_\pm=\log\pi_\pm\) & Real finite\\
\(\pi_\pm\) & Activity-premium multipliers & Strictly positive\\
\(a_\pm(\tilde v_h)\) & Side compensated unit-clock log moment & Complex transform domain\\
\(\eta_\pm\) & Endogenous return-clock exposure multiplier & Strictly positive\\
\(u_\pm^B,u_\pm^E\) & Affine baseline and endogenous activity loadings & Complex finite\\
\(z_\pm\) & Affine event-count multipliers & Strictly positive real for pricing\\
\(q_h\) & Affine news-return loading for interval \(h\) & \(q_h=\tilde v_h\)\\
\bottomrule
\end{tabularx}

\section{Assumptions used}

\begin{tabularx}{\textwidth}{XXX}
\toprule
Assumption & Where it is used & Why it is needed\\
\midrule
The terminal deterministic component is known at valuation time.
& Keeping \(d_{t,T}\) outside the stochastic cumulant.
& It contributes a scalar phase or log-moment shift, not an affine state loading.\\
Terminal mean reversion maps interval core increments linearly with weights \(w_h\).
& Replacing \(v\) by \(\tilde v_h=w_hv\) interval by interval.
& Linear weighting preserves transform composition across intervals.\\
Market-pressure anchor movement is signed-balance exposure, not a core shock.
& Passing \(v\) as the pressure return loading and \(\rho_\Omega\) as the pricing-kernel exposure loading.
& The pressure transform applies \(v\gamma_\Omega-\rho_\Omega\) to signed-balance exposure.\\
The conditional trading block has side-additive compensated log moments.
& Defining \(a_+(\tilde v_h)\) and \(a_-(\tilde v_h)\).
& The affine oracle consumes activity-clock loadings, not raw NIG parameters.\\
Baseline activity and endogenous excitation have separate return-clock exposures.
& Mapping \(a_\pm\) to \(u_\pm^B\) and \(u_\pm^E\).
& Pricing queries must preserve the difference between predictable baseline exposure and excitation-driven exposure.\\
Activity premia are count multipliers.
& Setting \(z_\pm=\pi_\pm\).
& The density tilts event counts directly through \(z_\pm^{N^\pm}\).\\
\bottomrule
\end{tabularx}

\section{Result}

\noindent\textbf{Identity.} For interval \(h\), set
\[
    \tilde v_h=w_h v.
\]
Let \(o_h\) be the materialized open-market exposure fraction for the interval.
Define the side compensated unit-clock log moments
\[
    a_+(\tilde v_h)
    =
    \log\mathbb E\!\left[e^{\tilde v_h(L_1^+-m_+)}\right],
    \qquad
    a_-(\tilde v_h)
    =
    \log\mathbb E\!\left[e^{-\tilde v_h(L_1^--m_-)}\right].
\]
Then the affine schedule query receives
\[
    \begin{aligned}
        u_{+,h}^B&=-a_+(\tilde v_h),
        &
        u_{+,h}^E&=-\eta_+a_+(\tilde v_h),
        &
        z_{+,h}&=\pi_+,
        \\
        u_{-,h}^B&=-a_-(\tilde v_h),
        &
        u_{-,h}^E&=-\eta_-a_-(\tilde v_h),
        &
        z_{-,h}&=\pi_-,
        \\
        r_{\Omega,h}&=v,
        &
        \rho_{\Omega,h}&=\rho_\Omega,
        &
        q_h&=\tilde v_h.
    \end{aligned}
\]
With these interval inputs, the affine schedule exponent equals
\[
    K_{t,T}(v,\log\pi_+,\log\pi_-,\rho_\Omega).
\]

\section{Derivation}

For one interval, the conditional trading-block MGF from
\texttt{03\_conditional\_cf\_mgf\_trading\_block.tex} is
\[
    \exp\!\left(
        A_h^+\left[K_+(\tilde v_h)-\tilde v_h m_+\right]
        +
        A_h^-\left[K_-(-\tilde v_h)+\tilde v_h m_-\right]
    \right).
\]
Write
\[
    a_+(\tilde v_h)=K_+(\tilde v_h)-\tilde v_h m_+,
    \qquad
    a_-(\tilde v_h)=K_-(-\tilde v_h)+\tilde v_h m_-.
\]
The effective return clocks are
\[
    A_h^+=B_{Z,h}^++\eta_+E_h^+,
    \qquad
    A_h^-=B_{Z,h}^-+\eta_-E_h^-.
\]
Thus, the trading contribution to the interval transform is
\[
    \exp\!\left(
a_+(\tilde v_h)B_{Z,h}^+
+\eta_+a_+(\tilde v_h)E_h^+
+a_-(\tilde v_h)B_{Z,h}^-
+\eta_-a_-(\tilde v_h)E_h^-
\right).
\]

The affine schedule transform uses the convention
\[
    \exp\!\left(
-u_+^B B_Z^+-u_+^E E^+
-u_-^B B_Z^--u_-^E E^-
+\omega\Omega_\Delta
+r_\Omega\int\Omega_s\,ds
\right).
\]
Matching exponents gives
\[
    u_{+,h}^B=-a_+(\tilde v_h),
    \qquad
    u_{+,h}^E=-\eta_+a_+(\tilde v_h),
    \qquad
    u_{-,h}^B=-a_-(\tilde v_h),
    \qquad
    u_{-,h}^E=-\eta_-a_-(\tilde v_h).
\]

The pricing density uses activity-premium factors
\[
    \pi_+^{N_{t,T}^+}\pi_-^{N_{t,T}^-}
    =
    \prod_h \pi_+^{N_h^+}\pi_-^{N_h^-}.
\]
Because the affine schedule transform represents count tilts as
\(z_+^{N_h^+}z_-^{N_h^-}\), the count multipliers are
\[
    z_{+,h}=\pi_+,
    \qquad
    z_{-,h}=\pi_-.
\]

Open news and reopening gaps are additive log-return components, so the same
weighted terminal argument applies to their return transform:
\[
    q_h=\tilde v_h.
\]

Market-pressure anchor movement is instead
\[
    \gamma_\Omega\int_{I_h} b_{\Omega_s}\,dE_s^\Omega.
\]
The pricing kernel additionally contributes
\[
    -\rho_\Omega\int_{I_h} b_{\Omega_s}\,dE_s^\Omega.
\]
Therefore the pressure transform applies the diagonal state loading
\[
    (v\gamma_\Omega-\rho_\Omega)e_h^\Omega b_k
\]
in pressure state \(k\). Its terminal return argument is not weighted by
\(w_h\). On the interval-grid representation this gives
\[
    r_{\Omega,h}=v,
    \qquad
    \rho_{\Omega,h}=\rho_\Omega,
\]
where \(e_h^\Omega\) carries the interval's signed-balance exposure scale.

The affine oracle composes intervals backward.
Since each interval receives the loadings above and the terminal interval
composition passes inherited \(\mathbf D\), \(V\), and pressure terminal values
to earlier intervals, the resulting schedule exponent is the stochastic
terminal cumulant
\[
    K_{t,T}(v,\log\pi_+,\log\pi_-,\rho_\Omega).
\]

\section{Domain and validity}

The adapter is valid when each weighted argument \(\tilde v_h=w_hv\) lies in
the side NIG moment domains needed for \(a_+\) and \(a_-\), the unweighted
market-pressure loading \(r_{\Omega,h}\) is finite, the induced affine
schedule query satisfies the domain constraints in
\texttt{06\_affine\_domain\_constraints.tex}, and the activity premia satisfy
\(\pi_+>0\) and \(\pi_->0\).
The pressure risk price \(\rho_\Omega\) must be finite, and the weights \(w_h\)
must be finite and non-negative.

\section{Consequences}

\begin{itemize}
    \item The terminal cumulant is computed by the same affine schedule oracle as physical activity transforms.
    \item Activity premia are count multipliers, not intensity compensators.
    \item Mean-reversion effects enter core shocks through interval weights \(w_h\) and the deterministic terminal component \(d_{t,T}\).
    \item Market-pressure anchor movement enters through signed-balance exposure and is not weighted by \(w_h\).
    \item Pressure risk enters as a pricing-kernel exposure loading, not as a second pressure generator.
    \item Baseline and endogenous loadings stay separate, preserving the return-clock exposure distinction needed by pricing.
\end{itemize}

\section{Checks}

\begin{tabularx}{\textwidth}{XXX}
\toprule
Check & Expected result & Evidence location\\
\midrule
Interval argument weighting
& The terminal argument is converted to a weighted interval argument before side cumulants and news loading are built.
& \texttt{build\_terminal\_affine\_schedule\_interval\_input} in \texttt{crates/janiform-engine-pricing/src/spectral/pricer/terminal\_affine/cumulant.rs}\\
Activity loading adapter
& Baseline loadings use \(-a_\pm\), endogenous loadings use \(-\eta_\pm a_\pm\), \(r_{\Omega,h}=v\), \(\rho_{\Omega,h}=\rho_\Omega\), and \(q_h=\tilde v_h\).
& \texttt{build\_terminal\_affine\_activity\_transform\_inputs} in \texttt{crates/janiform-engine-pricing/src/spectral/pricer/terminal\_affine/cumulant.rs}\\
Count-premium adapter
& Activity premia scale \(z_\pm\) without adding an intensity compensator.
& Test \texttt{terminal\_density\_activity\_premium\_tilts\_counts\_without\_intensity\_compensator}\\
Return weights
& Interval terminal weights equal the product of future mean-reversion decay factors for core shocks only.
& Test \texttt{terminal\_return\_weights\_apply\_future\_mean\_reversion\_decay}\\
\bottomrule
\end{tabularx}

\end{document}
